\documentclass[10pt]{letter}
\usepackage[margin=0.78in]{geometry}
\usepackage{microtype}
\signature{Abdulaziz M. D. Aljohani\\Corresponding author\\University of Tabuk, Saudi Arabia}
\address{Faculty of Computers and Information Technology\\University of Tabuk\\Tabuk 47512, Saudi Arabia}
\begin{document}
\begin{letter}{Guest Editors\\Special Issue: Cloud/Edge Computing for Next-Generation Networks\\\emph{Future Internet}}
\opening{Dear Guest Editors,}

We submit the manuscript entitled ``Guarded Multi-Fidelity Optimization with Trace-Derived Risk Auditing for Privacy-Constrained Cloud--Edge LLM Serving'' for consideration as an Article in the Special Issue ``Cloud/Edge Computing for Next-Generation Networks: Architecture and Applications.''

The paper studies delayed cloud activation for a hybrid LLM service in which privacy-sensitive requests remain local and eligible traffic may use remote capacity. The decision set contains 190 activation--release threshold pairs. DR-BGS uses a reflected diffusion only as a low-cost trend, while structural anchors, residual Gaussian-process learning, and a compound-jump simulator control the high-fidelity search. Across 36 exhaustively evaluated synthetic scenarios and five initializations, all 180 selected policies are within 1\% of exhaustive training selection on independent holdout streams, with 0.348\% maximum excess and 88.77\% fewer high-fidelity replications. A same-anchor GP reaches 0.878\%, showing that structural coverage supplies most of the observed reliability.

The trace-derived study separates efficient nomination from operational feasibility and statistical certification. Under strict protected- and eligible-blocking limits of 2\% and 5\%, all 120 method runs abstain. Exhaustive evaluation of all 2,280 policy--scenario combinations confirms that no policy in the tested family is certifiable at the original load ratios, and increasing certification evidence does not alter that result. A separately labeled operating-envelope analysis identifies substantial load reduction as the most consistent route to feasibility. At a point fixed before fresh result generation, $K=40$ and $\rho=0.05$, every unique DR-BGS, guarded-GP, and exhaustive-training nomination passes two disjoint 1,200-replay batches. The manuscript presents this as an operating-envelope diagnosis, not a deployment recommendation.

To address the limited basis of the original 95-positive-record trace slice, we added a separately locked robustness replication using the first seven days of the official \texttt{BurstGPT\_1.csv} release: 11,810 positive-token search records and 6,086 later certification records. The larger window reproduces both main findings. It yields 0/2,280 certifiable high-load policy--scenario pairs, while all 18 cross-method-distinct low-load scenario--policy pairs pass both certification batches. Source-file provenance, protocol locks, raw tables, code, and hashes are included in the release.

The contribution is an evidence-gated cloud--edge optimization workflow that distinguishes simulation-efficient nomination, operational feasibility, and independent certification. The paper does not claim production safety, full-corpus representativeness, DR-BGS superiority in every shifted regime, or a new general theorem for safe or multi-fidelity Bayesian optimization.

We confirm that neither the manuscript nor any part of its content is currently under consideration for publication in, or published in, another journal. All authors have approved the manuscript and agree with its submission to \emph{Future Internet}. The authors declare no conflicts of interest.

\closing{Sincerely,}
\end{letter}
\end{document}
